\section*{Hoofdstuk \ref{ch:b1:seq} --- Rijen}

\begin{solution}{exo:b1:seq:1}
$\Bigl|\dfrac{2n+1}{n+3} - 2\Bigr| = \dfrac{5}{n+3}$. Neem bij gegeven
$\varepsilon > 0$ een $N > \frac 5\varepsilon - 3$ (Archimedes): voor $n
\geq N$ is $\frac{5}{n+3} \leq \varepsilon$. De limiet is dus $2$.

$((-1)^n)$: haar \omterm{def:b1:seq:subsequence}{deelrijen} $(u_{2n}) = (1)$ en $(u_{2n+1}) = (-1)$
convergeren naar verschillende limieten, dus divergeert de rij
(\cref{prop:b1:seq:subsequences}). (Rechtstreeks: elke kandidaat $\ell$
faalt voor $\varepsilon = \frac12$, want opeenvolgende termen liggen op
afstand $2$.)
\end{solution}

\begin{solution}{exo:b1:seq:2}
Deling door $n^2$: $\dfrac{1 - 3/n + 1/n^2}{2 + 5/n^2} \to \dfrac12$.

$\sqrt{n+1} - \sqrt n = \dfrac{1}{\sqrt{n+1} + \sqrt n} \to 0$ (met de
toegevoegde uitdrukking).

$\dfrac{2^n + n^3}{3^n - n^2} = \dfrac{(2/3)^n + n^3/3^n}{1 - n^2/3^n}
\to \dfrac{0 + 0}{1 - 0} = 0$, met $q^n \to 0$ voor $\abs q < 1$ en de
vergelijking tussen \omterm{def:b1:poly:def}{veelterm} en meetkundige groei
(\cref{prop:b1:functions:powerrules}).

$5^{1/n} = \eu^{\frac{\ln 5}{n}} \to \eu^0 = 1$.
\end{solution}

\begin{solution}{exo:b1:seq:3}
Bernoulli: $(1+h)^n \geq 1 + nh$ voor $h \geq -1$, met inductie ---
$(1+h)^{n+1} = (1+h)^n(1+h) \geq (1+nh)(1+h) = 1 + (n+1)h + nh^2 \geq 1 +
(n+1)h$.

Voor $0 < \abs q < 1$: schrijf $\frac{1}{\abs q} = 1 + h$ met $h > 0$;
dan is $\abs{q}^n = \frac{1}{(1+h)^n} \leq \frac{1}{1 + nh} \to 0$, en de
insluiting geeft $q^n \to 0$ (het geval $q = 0$ is triviaal). Voor $q =
1$: een constante rij met limiet $1$. Voor $q = -1$: divergentie
(\cref{exo:b1:seq:1}). Voor $\abs q > 1$: $\abs q^n = (1 + h)^n \geq 1 +
nh \to +\infty$, zodat $(q^n)$ onbegrensd en dus divergent is (naar
$+\infty$ als $q > 1$; met wisselende tekens en zonder limiet als $q <
-1$).
\end{solution}

\begin{solution}{exo:b1:seq:4}
Vast punt: $\ell = \frac{\ell + 3}{2}$ geeft $\ell = 3$. Dan is
\[
v_{n+1} = u_{n+1} - 3 = \frac{u_n + 3}{2} - 3 = \frac{u_n - 3}{2} =
\frac{v_n}{2}:
\]
$(v_n)$ is meetkundig met reden $\frac12$ en $v_0 = -3$. Dus $u_n = 3 -
\frac{3}{2^n} \to 3$.
\end{solution}

\begin{solution}{exo:b1:seq:5}
$H_{2n} - H_n = \sum_{k=n+1}^{2n} \frac 1k \geq n \cdot \frac{1}{2n} =
\frac12$ (elk van de $n$ termen is $\geq \frac{1}{2n}$). Was $(H_n)$ een
\omterm{def:b1:seq:cauchy}{Cauchyrij}, dan zou $\varepsilon = \frac13$ nemen afdwingen dat
$\abs{H_{2n} - H_n} \leq \frac13$ voor grote $n$: tegenspraak. Een
stijgende niet-convergente rij divergeert naar $+\infty$
(\cref{thm:b1:seq:monotone}): $H_n \to +\infty$.
\end{solution}

\begin{solution}{exo:b1:seq:6}
Zij $a = \lim u_{2n}$, $b = \lim u_{2n+1}$ en $c = \lim u_{3n}$. De rij
$(u_{6n})$ is een \omterm{def:b1:seq:subsequence}{deelrij} van zowel $(u_{2n})$ als $(u_{3n})$: haar
limiet is gelijk aan $a$ en aan $c$, dus $a = c$. De rij $(u_{6n+3})$ is
een \omterm{def:b1:seq:subsequence}{deelrij} van $(u_{2n+1})$ (oneven indices) en van $(u_{3n})$ (indices
$6n + 3 = 3(2n+1)$): dus $b = c$. Bijgevolg is $a = b$, en
\cref{prop:b1:seq:subsequences} (even en oneven met gelijke limieten)
geeft de convergentie van $(u_n)$.
\end{solution}

\begin{solution}{exo:b1:seq:7}
\emph{Stabiliteit en grenzen:} $I = \intcc{0}{2}$ is stabiel: voor $x \in
I$ is $\sqrt{2 + x} \in \intcc{\sqrt 2}{2} \subseteq I$; en $u_0 = 0 \in
I$.

\emph{Monotonie:} $f(x) = \sqrt{2+x}$ is stijgend en $u_1 = \sqrt 2 >
u_0$: met inductie is $(u_n)$ stijgend. Stijgend en naar boven begrensd
door $2$: ze convergeert (\cref{thm:b1:seq:monotone}).

\emph{Limiet:} $\ell = \sqrt{2 + \ell}$ met $\ell \geq 0$ geeft $\ell^2 -
\ell - 2 = 0$, dus $\ell = 2$.

\emph{Foutafschatting:} vermenigvuldigen met de toegevoegde uitdrukking
geeft
\[
2 - u_{n+1} = 2 - \sqrt{2 + u_n}
= \frac{4 - (2 + u_n)}{2 + \sqrt{2+u_n}}
= \frac{2 - u_n}{2 + \sqrt{2 + u_n}}
\leq \frac{2 - u_n}{3},
\]
want $\sqrt{2 + u_n} \geq \sqrt 2 > 1$. Met inductie vanaf $2 - u_0 = 2$:
$\;0 \leq 2 - u_n \leq \frac{2}{3^n}$.
\end{solution}

\begin{solution}{exo:b1:seq:8}
Stel dat $(u_n)$ niet naar $\ell$ convergeert: voor een zekere
$\varepsilon_0 > 0$ voldoen oneindig veel indices aan $\abs{u_n - \ell} >
\varepsilon_0$; zij vormen een \omterm{def:b1:seq:subsequence}{deelrij} $(u_{\varphi(n)})$. Die \omterm{def:b1:seq:subsequence}{deelrij} is
begrensd, dus heeft ze volgens Bolzano--Weierstrass
(\cref{thm:b1:seq:bw}) een convergente \omterm{def:b1:seq:subsequence}{deelrij}, waarvan de limiet $\ell'$
voldoet aan $\abs{\ell' - \ell} \geq \varepsilon_0$ (de ongelijkheid gaat
mee naar de limiet, \cref{thm:b1:seq:order}). Maar een \omterm{def:b1:seq:subsequence}{deelrij} van een
\omterm{def:b1:seq:subsequence}{deelrij} van $(u_n)$ is een convergente \omterm{def:b1:seq:subsequence}{deelrij} van $(u_n)$, dus is
volgens de hypothese $\ell' = \ell$: tegenspraak.
\end{solution}

\begin{solution}{exo:b1:seq:9}
Stel $\eu = \frac pq$ met $q \geq 1$. De strikte ongelijkheden $a_q < \eu
< a_q + \frac{1}{q\,q!}$ (strikt omdat $(a_n)$ strikt stijgend en $(b_n)$
strikt dalend is), vermenigvuldigd met $q!$, geven
\[
q!\,a_q \;<\; q!\,\frac pq \;<\; q!\,a_q + \frac 1q \leq q!\,a_q + 1.
\]
Nu is $N = q!\,a_q = \sum_{k=0}^{q} \frac{q!}{k!}$ een geheel getal (elke
$\frac{q!}{k!}$ is voor $k \leq q$ een product van gehele getallen), en
$q!\,\frac pq = (q-1)!\,p$ eveneens. De formule legt het gehele getal
$(q-1)!\,p$ dus strikt tussen $N$ en $N + \frac 1q \leq N + 1$: een
geheel getal strikt binnen $\intoo{N}{N+1}$, wat onmogelijk is. Bijgevolg
is $\eu \notin \Q$.
\end{solution}

\begin{solution}{exo:b1:seq:10}
\begin{enumerate}
  \item Zij $\varepsilon > 0$ en $N$ met $\abs{u_k - \ell} \leq
        \frac{\varepsilon}{2}$ voor $k > N$. Voor $n > N$ is
        \[
        \abs{c_n - \ell}
        = \Bigl|\frac{\sum_{k=1}^{n}(u_k - \ell)}{n}\Bigr|
        \leq \frac{\sum_{k=1}^{N} \abs{u_k - \ell}}{n}
        + \frac{n - N}{n}\cdot\frac{\varepsilon}{2}
        \leq \frac{C}{n} + \frac{\varepsilon}{2},
        \]
        waarbij $C = \sum_{k=1}^N \abs{u_k - \ell}$ vast is. Voor grote
        $n$ is $\frac Cn \leq \frac\varepsilon2$: dan is $\abs{c_n -
        \ell} \leq \varepsilon$.
  \item $u_n = (-1)^n$: divergeert, en toch is $c_n \to 0$ (de partiële
        sommen zijn door $1$ begrensd, gedeeld door $n$).
  \item Pas (1) toe op de rij $v_n = u_{n+1} - u_n \to \ell$: haar
        gemiddelde van Cesàro is $\frac{u_{n+1} - u_1}{n} \to \ell$
        (telescoperen), en $\frac{u_{n+1}}{n} = \frac{u_{n+1} - u_1}{n}
        + \frac{u_1}{n} \to \ell$; de indices hernormaliseren
        ($\frac{u_n}{n} = \frac{u_n}{n-1}\cdot\frac{n-1}{n}$) geeft
        $\frac{u_n}{n} \to \ell$.
\end{enumerate}
\end{solution}

\begin{solution}{exo:b1:seq:11}
Zij $L = \inf_{n \geq 1} \frac{u_n}{n} \geq 0$ en $\varepsilon > 0$. Kies
$m$ met $\frac{u_m}{m} \leq L + \varepsilon$. Elke $n$ schrijft zich als
$n = qm + r$ met $0 \leq r < m$; de (herhaalde) subadditiviteit geeft
$u_n \leq q\,u_m + u_r$, dus
\[
\frac{u_n}{n} \leq \frac{q m}{n}\cdot\frac{u_m}{m} + \frac{u_r}{n}
\leq \frac{u_m}{m} + \frac{\max(u_0, \dots, u_{m-1})}{n}
\leq L + \varepsilon + \frac{C_m}{n},
\]
met $qm \leq n$. Voor grote $n$ is $\frac{C_m}{n} \leq \varepsilon$: dus
$L \leq \frac{u_n}{n} \leq L + 2\varepsilon$ voor alle grote $n$, en dat
is de convergentie naar $L$.
\end{solution}

\begin{solution}{exo:b1:seq:12}
De \omterm{def:b1:structures:subgroup}{deelgroep} $G = \Z + 2\pi\Z$ van $(\R, +)$ ligt dicht: ze is niet
$\alpha\Z$, want $1 = p\alpha$ en $2\pi = q\alpha$ zouden $2\pi = \frac
qp$ rationaal maken --- en $\pi \notin \Q$ (hier aangenomen; een bewijs
wordt geschetst in \cref{ch:b1:integration}). Volgens
\cref{exo:b1:reals:9} ligt $G$ dus dicht in $\R$.

Zij nu $y \in \intcc{-1}{1}$ en $\theta = \arcsin y$. Wegens de dichtheid
zijn er voor elke $\varepsilon > 0$ getallen $n \in \Z$ en $k \in \Z$ met
$\abs{(n + 2\pi k) - \theta} \leq \varepsilon$, dus ligt $n$ binnen
$\varepsilon$ van $\theta - 2\pi k$; omdat $\sin$ $2\pi$-periodiek en
$1$-lipschitz is ($\abs{\sin a - \sin b} \leq \abs{a - b}$, uit de
middelwaardeongelijkheid van \cref{ch:b1:derivative}), is dan
\[
\abs{\sin n - y} = \abs{\sin(n + 2\pi k) - \sin\theta}
\leq \abs{n + 2\pi k - \theta} \leq \varepsilon .
\]
Eén detail: $n$ doorloopt $\Z$, maar $\sin(-n) = -\sin n$ en $y$ was
willekeurig in $\intcc{-1}{1}$, zodat niet-negatieve indices volstaan
(vervang zo nodig $(n, y)$ door $(-n, -y)$). Bijgevolg ligt $\{\sin n : n
\in \N\}$ dicht in $\intcc{-1}{1}$; een rij die dicht ligt in een segment
heeft \omterm{def:b1:seq:subsequence}{deelrijen} die verschillende waarden naderen, en divergeert dus.
\end{solution}

\begin{solution}{pb:b1:seq:1}
\textbf{1.} $\dfrac{n(n+1)/2}{n^2} = \dfrac{1 + 1/n}{2} \to \dfrac12$, en
$\dfrac{n(n+1)(2n+1)/6}{n^3} = \dfrac{(1 + 1/n)(2 + 1/n)}{6} \to
\dfrac13$.

\textbf{2.} Boven: elk van de $n$ termen is $\leq \sqrt n$, dus $T_n \leq
n\sqrt n$. Onder: de termen met $k > \frac n2$ zijn er minstens $\frac
n2$, en elk is $\geq \sqrt{n/2}$:
\[
T_n \geq \frac n2 \sqrt{\frac n2} = \frac{n^{3/2}}{2\sqrt 2} .
\]

\textbf{3.} Voor $k \geq N$ geldt, omdat $b_{k+1} - b_k > 0$, dat
$m(b_{k+1} - b_k) \leq a_{k+1} - a_k \leq M(b_{k+1} - b_k)$. Sommeren
voor $k = N, \dots, n - 1$ laat beide leden telescoperen:
\[
m\,(b_n - b_N) \leq a_n - a_N \leq M\,(b_n - b_N),
\]
en deling door $b_n - b_N > 0$ geeft de bewering.

\textbf{4.} Per definitie van de limiet is er een $N$ met $\ell -
\varepsilon \leq \frac{a_{k+1} - a_k}{b_{k+1} - b_k} \leq \ell +
\varepsilon$ voor alle $k \geq N$; vraag 3 met $m = \ell - \varepsilon$
en $M = \ell + \varepsilon$ draagt de insluiting over op $\frac{a_n -
a_N}{b_n - b_N}$.

\textbf{5.} Het rechterlid van de identiteit uitwerken:
\[
\frac{a_N - \ell b_N}{b_n} + \frac{a_n - a_N}{b_n} -
\ell\,\frac{b_n - b_N}{b_n}
= \frac{a_n - \ell b_n}{b_n} = \frac{a_n}{b_n} - \ell .
\]
Volgens vraag 4 is de tweede factor van het product in absolute waarde
door $\varepsilon$ begrensd, en $0 < 1 - \frac{b_N}{b_n} \leq 1$ voor
grote $n$, dus
\[
\Bigl|\frac{a_n}{b_n} - \ell\Bigr|
\leq \frac{\abs{a_N - \ell b_N}}{b_n} + \varepsilon
\leq 2\varepsilon
\]
zodra $b_n \geq \frac{\abs{a_N - \ell b_N}}{\varepsilon}$, wat
uiteindelijk gebeurt omdat $b_n \to +\infty$. Bijgevolg is
$\frac{a_n}{b_n} \to \ell$: de stelling van Cesàro--Stolz.

\textbf{6.} Kies bij gegeven $M$ een $N$ met $\frac{a_{k+1} -
a_k}{b_{k+1} - b_k} \geq M$ voor $k \geq N$; de onderste helft van vraag
3 geeft $a_n - a_N \geq M(b_n - b_N)$, dus
\[
\frac{a_n}{b_n} \geq \frac{a_N}{b_n} +
M\Bigl(1 - \frac{b_N}{b_n}\Bigr) \longrightarrow M .
\]
Voorbij een zekere rang is $\frac{a_n}{b_n} \geq M - 1$; omdat $M$
willekeurig was, gaat $\frac{a_n}{b_n} \to +\infty$.

\textbf{7.} Met $b_n = n$ en $a_n = u_1 + \dots + u_n$ is het quotiënt
van de aangroeiingen $u_{n+1} \to \ell$, zodat het gemiddelde van Cesàro
$\frac{a_n}{n}$ naar $\ell$ gaat: onderdeel (1) van
\cref{exo:b1:seq:10}. Met $a_n = u_n$ is het quotiënt van de
aangroeiingen $u_{n+1} - u_n$, wat onderdeel (3) geeft. Omkering: $a_n =
(-1)^n$ met $b_n = n$ heeft $\frac{a_n}{b_n} \to 0$, en toch wisselt
$a_{n+1} - a_n = \pm 2$: het quotiënt van de aangroeiingen heeft geen
limiet.

\textbf{8.} Toegevoegde uitdrukking:
\[
\bigl((1+h)^{3/2} - 1\bigr)\bigl((1+h)^{3/2} + 1\bigr)
= (1+h)^3 - 1 = 3h + 3h^2 + h^3 .
\]
Voor $h = \frac1n$: $n\bigl((1 + \frac1n)^{3/2} - 1\bigr) = \frac{3 + 3/n
+ 1/n^2}{(1 + 1/n)^{3/2} + 1}$, en $1 \leq (1 + \frac1n)^{3/2} \leq (1 +
\frac1n)^2 \to 1$ (insluiting), dus gaat de noemer naar $2$ en de
uitdrukking naar $\frac32$. Nu Stolz met $a_n = T_n$ en $b_n = n^{3/2}$
(strikt stijgend, $\to +\infty$):
\[
\frac{\sqrt{n+1}}{(n+1)^{3/2} - n^{3/2}}
= \frac{\sqrt{n+1}}{\sqrt n} \cdot
\frac{1}{n\bigl((1 + \frac1n)^{3/2} - 1\bigr)}
\longrightarrow 1 \cdot \frac{2}{3},
\]
dus $T_n \sim \frac23\,n^{3/2}$. (De insluiting van vraag 2 had de
constante in $\intcc{0.35}{1}$ opgesloten; Stolz pint haar vast.)

\textbf{9.} (G1) in $u = \ln(1+t)$: $1 + t = \eu^{\ln(1+t)} \geq 1 +
\ln(1+t)$, dus $\ln(1+t) \leq t$. (G1) in $u = -\frac{t}{1+t}$:
$\eu^{-t/(1+t)} \geq 1 - \frac{t}{1+t} = \frac{1}{1+t} > 0$; $\ln$ nemen
(stijgend) geeft $-\frac{t}{1+t} \geq -\ln(1+t)$, oftewel $\ln(1+t) \geq
\frac{t}{1+t}$.

\textbf{10.} $\ln$ is strikt stijgend (\cref{prop:b1:functions:expln}) en
$\ln(2^k) = k\ln 2$ is onbegrensd, dus $\ln n \to +\infty$.
Aangroeiingen: met $t = \frac1n$ in vraag 9 is
\[
\frac 1{n+1} = \frac{1/n}{1 + 1/n} \leq
\ln\Bigl(1 + \frac1n\Bigr) \leq \frac1n
\quad\Longrightarrow\quad
\frac{n}{n+1} \leq \frac{1/(n+1)}{\ln(1 + 1/n)} \leq 1 ,
\]
zodat het quotiënt van de aangroeiingen $\frac{H_{n+1} -
H_n}{\ln(n+1) - \ln n}$ naar $1$ gaat; Stolz geeft $H_n \sim \ln n$.

\textbf{11.} Zet $x_n = \frac{u_{n+1}}{u_n} \to L$ en $t_n =
\frac{x_n}{L} - 1 \to 0$. Vraag 9: $\frac{t_n}{1 + t_n} \leq \ln(1 + t_n)
\leq t_n$, dus $\ln x_n - \ln L = \ln(1 + t_n) \to 0$ door insluiting.
Cesàro (vraag 7) toegepast op $(\ln x_k)$:
\[
\frac1n \sum_{k=0}^{n-1} \ln x_k = \frac{\ln u_n - \ln u_0}{n}
\longrightarrow \ln L ,
\]
dus $\frac{\ln u_n}{n} \to \ln L$. Met $h_n = \frac{\ln u_n}{n} - \ln L
\to 0$ is $u_n^{1/n} = L\,\eu^{h_n}$, en (G1) sluit $1 + h_n \leq
\eu^{h_n} \leq \frac{1}{1 - h_n}$ in (voor $h_n < 1$), zodat $\eu^{h_n}
\to 1$ en $u_n^{1/n} \to L$. Toepassing: $u_n = \binom{2n}{n}$ geeft
\[
\frac{u_{n+1}}{u_n} = \frac{(2n+1)(2n+2)}{(n+1)^2}
= \frac{2(2n+1)}{n+1} \longrightarrow 4 ,
\qquad\text{dus}\qquad
\binom{2n}{n}^{1/n} \to 4 .
\]

\textbf{12.} Voor $0 < x \leq 1$ geeft (G2) dat $\sin x \geq x(1 -
\frac{x^2}{6}) \geq \frac{5x}{6} > 0$ en
\[
x - \sin x \geq \frac{x^3}{6} - \frac{x^5}{120}
= x^3\Bigl(\frac16 - \frac{x^2}{120}\Bigr)
\geq \frac{19}{120}\,x^3 > 0 :
\]
dus $0 < \sin x < x$ op $\intoc{0}{1}$. Steeds is $u_1 = \sin u_0 \in
\intcc{-1}{1}$. Is $u_1 = 0$, dan is $u_n = 0$ voor $n \geq 1$. Is $u_1
\in \intoc{0}{1}$, dan is met inductie $0 < u_{n+1} = \sin u_n < u_n \leq
1$, zodat $(u_n)_{n\geq1}$ strikt dalend en naar onderen begrensd door
$0$ is: ze convergeert naar een zekere $\ell \in \intco{0}{1}$
(\cref{thm:b1:seq:monotone}). De formule van product naar som en (G2)
geven $\abs{\sin a - \sin b} = 2\abs{\cos \frac{a+b}{2}}\,
\abs{\sin\frac{a-b}{2}} \leq \abs{a - b}$, dus $u_{n+1} = \sin u_n \to
\sin \ell$: $\ell = \sin\ell$. Was $\ell > 0$, dan zou $\sin\ell < \ell$
zijn: onmogelijk. Dus $u_n \to 0$.

\textbf{13.} Deling van (G2) door $u_n > 0$ geeft
\[
1 - \frac{u_n^2}{6} \leq \frac{\sin u_n}{u_n} \leq 1 -
\frac{u_n^2}{6} + \frac{u_n^4}{120} \leq 1 ,
\]
en $u_n \to 0$ sluit $\frac{\sin u_n}{u_n} \to 1$ in. Deling van $x -
\sin x$ door $x^3$ geeft
\[
\frac16 - \frac{u_n^2}{120} \leq \frac{u_n - \sin u_n}{u_n^3}
\leq \frac16 \longrightarrow \frac16 .
\]

\textbf{14.} Omdat $u_{n+1} = \sin u_n$, is
\[
w_n = \frac{u_n^2 - \sin^2 u_n}{u_n^2 \sin^2 u_n}
= \frac{(u_n - \sin u_n)(u_n + \sin u_n)}{u_n^2 \sin^2 u_n}
= \frac{u_n - \sin u_n}{u_n^3}\cdot
\frac{u_n + \sin u_n}{u_n}\cdot
\Bigl(\frac{u_n}{\sin u_n}\Bigr)^{2}
\]
(controleer de machten van $u_n$: $3 + 1 + (-4)$ tegenover de $u_n^2$
beneden en de $u_n^4$ boven). Volgens vraag 13 gaan de drie factoren naar
$\frac16$, $2$ en $1$: $w_n \to \frac13$.

\textbf{15.} $v_n = \frac{1}{u_n^2}$ heeft aangroeiingen $v_{n+1} - v_n =
w_n \to \frac13$, dus $\frac{v_n}{n} \to \frac13$ volgens
\cref{exo:b1:seq:10}~(3): $n u_n^2 \to 3$. Vervolgens is
\[
\abs{\sqrt n\,u_n - \sqrt 3}
= \frac{\abs{n u_n^2 - 3}}{\sqrt n\,u_n + \sqrt 3}
\leq \frac{\abs{n u_n^2 - 3}}{\sqrt 3} \longrightarrow 0 :
\]
$\sqrt n\,u_n \to \sqrt 3$, oftewel $u_n \sim \sqrt{3/n}$.

\textbf{16.} Omdat $n u_n^2 \to 3$, is uiteindelijk $2 \leq n u_n^2 \leq
4$, dus $\sqrt{2/n} \leq u_n \leq 2/\sqrt n$. Is $u_n \leq 10^{-2}$ met
$n$ in dat bereik, dan is $2/n \leq 10^{-4}$: $n \geq 20\,000$; en
$\sqrt{3/n} = 10^{-2}$ bij $n = 30\,000$. De \omterm{ex:b1:seq:heron}{methode van Heron}
kwadrateert de fout bij elke stap --- het aantal cijfers verdubbelt ---
omdat in haar vaste punt de relevante richtingscoëfficiënt in \omterm{def:b1:complex:field}{modulus} $<
1$ is (de iteratie is er contraherend). Hier is $\sin' 0 = \cos 0 = 1$:
het vaste punt is neutraal, er bestaat geen meetkundige contractie, en de
afname wordt geregeld door de eerste niet-lineaire term $-\frac{x^3}{6}$
en is dus polynomiaal. Eén stap van Heron wint meer nauwkeurigheid dan
tienduizend stappen van de sinus.

\textbf{17.} Voor willekeurige $u_0$ is $u_1 = \sin u_0 \in
\intcc{-1}{1}$. Is $u_1 = 0$, dan is de rij nul vanaf rang $1$. Is $u_1 >
0$, dan geldt Deel IV letterlijk. Is $u_1 < 0$, zet dan $v_n = -u_n$: de
oneven pariteit van $\sin$ geeft $v_{n+1} = -\sin u_n = \sin(-u_n) = \sin
v_n$ met $v_1 \in \intoc{0}{1}$, dus $v_n \sim \sqrt{3/n}$, oftewel $u_n
\sim -\sqrt{3/n}$. In alle gevallen is $\abs{u_n} \sim \sqrt{3/n}$ (of is
de rij uiteindelijk $0$): de val is universeel, alleen het teken
herinnert zich het beginpunt.

\textbf{18.} Eerst is $\frac{u_{n+1}}{u_n} = 1 - \frac{u_n -
u_{n+1}}{u_n^2}\,u_n \to 1 - a \cdot 0 = 1$. Dan is
\[
\frac{1}{u_{n+1}} - \frac{1}{u_n}
= \frac{u_n - u_{n+1}}{u_n u_{n+1}}
= \frac{u_n - u_{n+1}}{u_n^2}\cdot\frac{u_n}{u_{n+1}}
\longrightarrow a \cdot 1 = a ,
\]
en \cref{exo:b1:seq:10}~(3) geeft $\frac{1}{n u_n} \to a$, oftewel $n u_n
\to \frac1a$.

\textbf{19.} $v_n = \frac{1}{u_n}$: $v_{n+1} = \frac{1 + u_n}{u_n} = v_n
+ 1$, dus $v_n = v_0 + n$ en
\[
u_n = \frac{u_0}{1 + n u_0} , \qquad
n u_n = \frac{n u_0}{1 + n u_0} \longrightarrow 1 .
\]
Controle van het lemma: $u_n - u_{n+1} = \frac{u_n^2}{1 + u_n}$, dus
$\frac{u_n - u_{n+1}}{u_n^2} = \frac{1}{1 + u_n} \to 1 = a$, en vraag 18
voorspelt $n u_n \to 1$: exacte overeenstemming.

\textbf{20.} Positiviteit met inductie ($\eu^{-u} > 0$); dalend omdat
$\eu^{-u_n} < 1$ voor $u_n > 0$; dus $u_n \to \ell \geq 0$
(\cref{thm:b1:seq:monotone}). Brug via de continuïteit: met $h_n = \ell -
u_n \to 0$ is $\eu^{-u_n} = \eu^{-\ell} \eu^{h_n} \to \eu^{-\ell}$ wegens
de insluiting $1 + h_n \leq \eu^{h_n} \leq \frac{1}{1 - h_n}$ uit (G1);
dus $\ell = \ell\, \eu^{-\ell}$, en $\ell > 0$ zou $\eu^{-\ell} = 1$
afdwingen, wat onwaar is: $\ell = 0$. Voor $t > 0$ geeft (G1) dat
$\eu^{-t} \geq 1 - t$ en $\eu^{-t} \leq \frac{1}{1 + t}$, dus
\[
\frac{1}{1 + t} \leq \frac{1 - \eu^{-t}}{t} \leq 1 .
\]
Met $t = u_n$: $\frac{u_n - u_{n+1}}{u_n^2} = \frac{1 -
\eu^{-u_n}}{u_n} \to 1$. Vraag 18 met $a = 1$ geeft $n u_n \to 1$, dus
$u_n \sim \frac1n$.

\textbf{21.} Net als in vraag 18 is $\frac{u_{n+1}}{u_n} = 1 - \frac{u_n
- u_{n+1}}{u_n^3}\,u_n^2 \to 1$. Dan is
\[
\frac{1}{u_{n+1}^2} - \frac{1}{u_n^2}
= \frac{(u_n - u_{n+1})(u_n + u_{n+1})}{u_n^2 u_{n+1}^2}
= \frac{u_n - u_{n+1}}{u_n^3}\cdot
\frac{u_n + u_{n+1}}{u_n}\cdot
\Bigl(\frac{u_n}{u_{n+1}}\Bigr)^2
\longrightarrow a \cdot 2 \cdot 1 = 2a ,
\]
en \cref{exo:b1:seq:10}~(3) geeft $\frac{1}{n u_n^2} \to 2a$: $n u_n^2
\to \frac{1}{2a}$. Voor de sinus is $a = \frac16$ (vraag 13): $n u_n^2
\to 3$, precies Deel IV.

\textbf{22.} Aangroeiingen: $a_{n+1} - a_n = 2^{-n} - 2^{-n-1} =
2^{-n-1} = b_{n+1} - b_n$, dus is het quotiënt van de aangroeiingen
constant $1$. En toch is $\frac{a_n}{b_n} = \frac{2 - 2^{-n}}{1 - 2^{-n}}
\to 2 \neq 1$. Het bewijs van vraag 5 begeeft het bij de randterm:
$\frac{a_N - \ell b_N}{b_n} \to 0$ vergde $b_n \to +\infty$; hier is (met
$\ell = 1$) $a_N - b_N = 1$ en $b_n \to 1$, zodat die term naar $1$ gaat
--- precies het resterende verschil $2 - 1$.

\textbf{23.} Stolz met $A_n = \sum_{k=1}^n H_k$ en $B_n = n\ln n$:
$B_{n+1} - B_n = \ln(n+1) + n\ln(1 + \frac1n) > 0$ en $B_n \to +\infty$.
Volgens vraag 9 is $\frac{n}{n+1} \leq n\ln(1 + \frac1n) \leq 1$, dus
$B_{n+1} - B_n = \ln(n+1) + \theta_n$ met $\frac12 \leq \theta_n \leq 1$.
Bijgevolg is
\[
\frac{A_{n+1} - A_n}{B_{n+1} - B_n}
= \frac{H_{n+1}}{\ln(n+1)}\cdot
\frac{1}{1 + \theta_n/\ln(n+1)} \longrightarrow 1 \cdot 1 = 1
\]
(vraag 10 voor de eerste factor; $\theta_n$ begrensd en $\ln(n+1) \to
\infty$ voor de tweede). Stolz besluit: $\sum_{k=1}^n H_k \sim n\ln n$.

\textbf{24.} Is $u_n \to \ell > 0$, dan gaat, net als in vraag 11, $\ln
u_n \to \ln\ell$ (de insluiting van vraag 9 op $\ln\frac{u_n}{\ell}$),
zodat de gemiddelden van Cesàro $\frac1n\sum_{k=1}^n \ln u_k \to
\ln\ell$, en de exponentiële brug van vraag 11 geeft $(u_1\cdots
u_n)^{1/n} = \exp\bigl(\frac1n\sum\ln u_k\bigr) \to \ell$. Is $u_n \to
+\infty$, dan is voor elke $M$ uiteindelijk $u_n \geq \eu^M$, dus $\ln
u_n \geq M$: $\ln u_n \to +\infty$; de versie van Cesàro met $+\infty$
(vraag 6 met $b_n = n$) geeft $\frac1n\sum \ln u_k \to +\infty$, en (G1)
($\eu^s \geq 1 + s$) stuurt het meetkundig gemiddelde naar $+\infty$. Met
$u_n = n$: $(n!)^{1/n} \to +\infty$.

\textbf{25.} (i) De volledigheid kwam alleen via de stelling van de
monotone limiet binnen, om de limieten in de vragen 12 en 20 voort te
brengen; de stelling van Cesàro--Stolz zelf is zuiver
$\varepsilon$-beheer en geldt ook over $\Q$. (ii) Stolz vervangt $\lim
\frac{a_n}{b_n}$ door de limiet van het quotiënt van de aangroeiingen,
precies zoals l'Hospital $\lim\frac fg$ door $\lim\frac{f'}{g'}$ vervangt
--- de differentiële tweelingzus rust op de middelwaardestelling van
\cref{ch:b1:derivative}. (iii) Vuistregel: is $f(x) = x - a\,x^{p+1} +
o(x^{p+1})$ in het neutrale vaste punt $0$, dan is $\frac{1}{u_{n+1}^p} -
\frac{1}{u_n^p} \to pa$ en $u_n \sim (pan)^{-1/p}$: contact van orde $p +
1$ levert afname $n^{-1/p}$ --- hoe vlakker de grafiek tegen de diagonaal
ligt, hoe trager de val. (iv) De constante: (G2) levert de
derdegraadscoëfficiënt $\frac16$; de ontbinding van vraag 14 verdubbelt
die tot de aangroeiing $\frac13$ van de telescoop; Cesàro maakt van
$\frac{1}{u_n^2}$ het getal $\frac n3$; omkeren en worteltrekken levert
$\sqrt{3/n}$.
\end{solution}
